\section{Limitation and null result}
\label{sec:limitations}

This section reports one limitation and one null result of the study.

\begin{itemize}
    \item \textbf{Single model.} All experiments use Qwen3-4B-Instruct-2507. Replicating across model families would require substantially more inference compute and time, this is left for future work.

    \item \textbf{Repair geometry is a null result by data starvation.} A companion question asked whether the hidden state shift from a failing attempt to its repair attempt carries a geometric signature of repair success. The data cannot support it. Of the 266 tasks whose first attempt fails, only 20 ever recover within the 5-attempt budget, 14 of them at the first repair, and the conditional pass rate among still-unsolved tasks collapses from 40.1\% at attempt 0 to 5.3\%, 2.0\%, 0.4\% and 0\% at attempts 1 through 4. Prior reports find that iterative self-repair plateaus after the first one or two attempts \citep{arimbur2026howmanytries}, and the decay here is even sharper, with the pass rate collapsing at the very first repair. No contrastive direction can be meaningfully estimated from 14 successful repairs in a 2560-dimensional space, so no claim about repair geometry is made in either direction, and the distinction between ``the signal does not exist'' and ``the data were not sufficient to detect it'' cannot be resolved with the present dataset.
\end{itemize}
